\chapter{Introduction}\label{ch:intro}

\section{Motivation}

Sixth-generation wireless networks must provide reliable connectivity and
timely awareness of the surrounding environment for applications such as
industrial automation, intelligent transport, and autonomous aerial systems
\autocite{dongSensingService6G2023,kaushikIntegratedSensingCommunications2024}.
\autocite{liuIntegratedSensingCommunications2022,liuJointRadarCommunication2020}
Integrated sensing and communication (ISAC) addresses this requirement by
sharing spectrum, radio hardware, and spatial degrees of freedom between data
delivery and sensing. The resulting benefit is compelling, but the shared
resources also create a fundamental allocation problem: power assigned to
communication, dedicated sensing, and interference protection cannot be chosen
independently.

Cell-free massive MIMO makes this problem both more promising and more
challenging. A central processing unit can coordinate distributed access points
(APs) to serve users cooperatively, thereby creating spatial diversity for
communication and sensing without fixed cell boundaries
\autocite{femeniasScalableCellfreeMassive2025,demirhanCellfreeISACMIMO2025}.
\autocite{maoCommunicationSensingRegion2024,liNetworkedISAC2024}
However, each AP has a finite transmit-power budget. Selecting which APs may
illuminate each target therefore changes the sensing geometry and the remaining
power available for globally cooperative communication. The decision must also
remain reliable when the communication channel state information (CSI) is
imperfect.

\section{Research Gap, Problem, and Contributions}

Existing cooperative ISAC formulations establish the value of spatial
cooperation, yet three connected gaps remain for the design considered here.
First, a sensing association is often difficult to interpret when the model
does not distinguish a target's dedicated sensing waveform from the globally
cooperative communication signal. Second, perfect-CSI formulations can lose
their communication quality-of-service (QoS) guarantee under channel error.
Third, a relaxed semidefinite or continuous association solution is not itself
a physical beamforming and topology design; rank-one and binary recovery must
be assessed after the relaxation is solved.

This thesis addresses these gaps through robust downlink resource allocation
for a cell-free ISAC network. The objective is to minimize total transmit power
while satisfying worst-case communication SINR, target detection, and tracking
accuracy requirements. The central design choice is a known,
target-specific sensing covariance for each target. Its binary AP--target
association variable authorizes dedicated sensing power only; it does not gate
the globally cooperative communication beamformers. This separation makes the
topology decision physically interpretable and permits each target's Fisher
information and PCRB requirement to be evaluated from its own probing
covariance.

\section{Thesis Organization}

Chapter~\ref{ch:literature_review} narrows the literature from cooperative
ISAC to the four design requirements that motivate this thesis: cooperative
resource allocation, robust QoS, multi-target tracking accuracy, and
recoverable sensing topology. Chapter~\ref{ch:system_model} then defines the
network roles, the robust communication model, and the target-specific sensing
model. Chapter~\ref{ch:solution_method} formulates the power-minimization
problem and derives the lifted LMI and dual-DC recovery framework. Chapter
\ref{ch:results} evaluates the proposed workflow through topology, robustness,
recovery, and scalability evidence. Finally,
Chapter~\ref{ch:future_work} answers the research problem from the accumulated
evidence and identifies the validation work that remains.
